\documentclass{article}
\usepackage[utf8]{inputenc}
\usepackage{amsmath}
\usepackage{geometry}
\geometry{margin=2cm}
\begin{document}

\section*{Questão 122 — ENEM 2024 — Ciências da Natureza}

\subsection*{Tema: Metabolismo Procarionte nos Ciclos Biogeoquímicos}

Esta questão aborda a ação de organismos procariontes nos ciclos do nitrogênio e do carbono, focando em transformações químicas do tipo redox que esses organismos realizam.

\subsection*{Interpretação do Enunciado e Estratégia de Resolução}
O enunciado solicita identificar em qual das etapas numeradas ocorre uma transformação redox semelhante à realizada pelos procariontes nitrificantes, ou seja, uma reação química com ganho ou perda de elétrons correspondente a oxidação da amônia.

Para responder, deve-se analisar o esquema que mostra etapas distintas dos ciclos do nitrogênio e do carbono, reconhecendo os processos metabólicos dos procariontes, especialmente as reações redox de oxidação no ciclo do nitrogênio, e comparar com outras transformações para encontrar a mais similar.

\subsection*{Análise dos Conceitos Envolvidos}
\begin{tabular}{|p{4cm}|p{6cm}|p{6cm}|}
\hline
\textbf{Conceito} & \textbf{Ideia-Chave} & \textbf{Aplicação na questão} \\ \hline
Nitrificação & Oxidação de amônia em nitrito e nitrato por procariontes nitrificantes & Explica transformação redox oxidativa no ciclo do nitrogênio vinculada à nitrificação (etapas do esquema). \\ \hline
Desnitrificação & Redução de nitratos em nitrogênio gasoso por procariontes desnitrificantes & Contrasta com nitrificação por ser reação redox de redução e não oxidação \\ \hline
Decomposição no ciclo do carbono & Oxidação do carbono orgânico em CO$_2$ por procariontes decompositores & Transformação redox oxidativa similar à nitrificação, representada na etapa 5 do esquema \\ \hline
\end{tabular}

\subsection*{Fundamentação Teórica}
\begin{itemize}
  \item Os ciclos biogeoquímicos do nitrogênio e do carbono dependem da ação de procariontes que realizam reações químicas essenciais para o movimento e transformação desses elementos.
  \item No ciclo do nitrogênio, os procariontes nitrificantes realizam oxidação de amônia para nitrito e nitrato, processos que envolvem processos redox — em particular oxidações.
  \item No ciclo do carbono, procariontes decompositores oxidam compostos orgânicos retornando CO$_2$ à atmosfera, uma transformação redox de oxidação similar à nitrificação, apesar de ocorrer em outro ciclo.
\end{itemize}

\subsection*{Resolução Detalhada}
Observando o esquema, temos:
\begin{enumerate}
  \item Etapa 1: Fixação do nitrogênio, conversão de N$_2$ em amônia, realizada por procariontes fixadores, não caracterizada por ser oxidação redox como a nitrificação.
  \item Etapa 2: Desnitrificação, redução de nitratos em N$_2$, reação redox de redução, não oxidação.
  \item Etapa 3: Fixação de CO$_2$ por procariontes fotossintetizantes, não caracterizado como transformação redox de oxidação similar à nitrificação.
  \item Etapa 4: Fixação de CO$_2$ por plantas; plantas não são procariontes e não realizam transformação redox similar a nitrificação de procariontes.
  \item Etapa 5: Oxidação do carbono orgânico fixado em CO$_2$ por procariontes decompositores, processo redox de oxidação semelhante à nitrificação.
\end{enumerate}

Assim, a etapa 5 é a que melhor se enquadra na descrição do enunciado, sendo a resposta correta.

\subsection*{Pulo do Gato}
Para identificar transformações redox semelhantes, deve-se focar no tipo da reação (oxidação ou redução) e no metabolismo específico dos procariontes envolvidos, diferenciando-os claramente de processos vegetais ou outros ciclos que não realizam oxidação similar.

\subsection*{Armadilhas para Evitar}
\begin{itemize}
  \item Confundir fixação de nitrogênio com nitrificação.
  \item Confundir desnitrificação (redução) com nitrificação (oxidação).
  \item Associar fixação de CO$_2$ por procariontes fotossintetizantes à oxidação redox nitrificante.
  \item Atribuir transformações metabólicas das plantas a procariontes.
\end{itemize}

\subsection*{Código da Questão}
CN\_CIENCIA\_VIDA\_001


\end{document}